\section{Varian Implikasi (Konvers, Invers) dan Negasi Majemuk}

Sebagaimana layaknya kalimat bersyarat di ranah hukum perundang-undangan, sebuah klausa Implikasi $p \rightarrow q$ ("Jika $p$, maka $q$") memiliki tiga buah bayangan atau 'varian' kalimat modifikasi yang kerang kali disalahgunakan oleh pengacara untuk membangun kesesatan berpikir (Fallacy). Membangun benteng pemahaman forensik di ketiga modifikasi ini mencegah \textit{Software Engineer} terjebak menulis kode persyaratan eror bersarang (*Logic Bug*) \cite{rosen2019}.

Ketiga varian mutasi dari panah Implikasi orisinal $p \rightarrow q$ dinamakan: Konvers, Invers, dan Kontraposisi.

\subsection{1. Pembalikan Arah Panah: Konvers ($q \rightarrow p$)}

Varian pertama lahir secara sembarangan mubah dengan cara menukar paksa kursi kemudi. Sang akibat $q$ dicaplok ditarik menduduki singgasana sang pemicu, sementara mantan pemicu asli didorong jatuh menjadi akibat semu.

\begin{contoh}[Analisis Kegagalan Logika Serampangan Konvers]
\textbf{Implikasi Asli Valid ($p \rightarrow q$):} "Jika Anda sukses bermukim ber-KTP di kota Depok ($p$), pastilah letak lokasi Anda tergaransi mutlak menancap di bentang provinsi Jawa Barat ($q$)." (Murni Valid Tautologi Semesta).

\textbf{Varian Konvers ($q \rightarrow p$):} "Jika lokasi Anda tengah melintasi bentang provinsi Jawa Barat ($q$), niscaya mutlak Anda tengah ber KTP mukim mendiami kota Depok ($p$)!!"

Apakah Varian Konvers ini Ekuivalen menggaransi kebenaran serupa dengan aslinya tadi? \textbf{TENTU TIDAK!} Seseorang di Jawa Barat bisa saja sedang pelesir makan pisang goreng di terminal Bogor atau berendam asyik di kubangan Bandung. 
\textbf{Kesimpulan Mutlak:} Konvers ($q \rightarrow p$) **TIDAK PERNAH** ekuivalen menjiplak karakter sejati saudara aslinya ($p \rightarrow q$). \cite{wiki_first_order_logic}
\end{contoh}

\subsection{2. Tangan Pengingkar Murni: Invers ($\neg p \rightarrow \neg q$)}

Alih-alih menyentuh posisi kursi mereka duduk, Invers menghajar lurus merusak sifat nilai kasta mereka masing-masing menjadi lawan sengketa dari kodrat asalnya. Ia menyuntikkan bisa Negasi ke kedua ruas serentak.

\begin{contoh}[Invers Menularkan Sesat Pikir]
\textbf{Implikasi Asli Valid ($p \rightarrow q$):} "Jika langit Jakarta memuntahkan galon curah Hujan keras sore ini ($p$), aspal jalanan Thamrin niscaya mekar berkubang Basah ($q$)." (Sah Benar Secara Fisika).

\textbf{Varian Invers ($\neg p \rightarrow \neg q$):} "Jika kenyataan awan melenggang permai tak memuntahkan Hujan detik ini ($\neg p$), niscaya sakti aspal Thamrin takkan pernah bisa Basah menyerap air setetespun ($\neg q$)!" 

Apakah pernyataan Invers naif ini Ekuivalen dengan fakta asalnya? \textbf{CACAT TOTAL!} Walau langit tidak menangis hujan meratapi nasib ($\neg p$), aspal jalan Thamrin bisa asyik basah kuyup tersiram mobil petugas Pemadam Kebakaran DKI yang bocor mancur tangkinya! 
\textbf{Kesimpulan Mutlak:} Invers ($\neg p \rightarrow \neg q$) **MUSTAHIL** memegang stempel ekuivalensi identik dengan asalnya ($p \rightarrow q$). \cite{wiki_first_order_logic}
\end{contoh}

\subsection{3. Sang Penyelamat Takhta Ekuivalen: Kontraposisi ($\neg q \rightarrow \neg p$)}

Inilah mahkota penyelamat para logikawan, senjata dekonstruksi rahasia terbaik saat menjumpai teorem buntu di pengadilan! Varian Kontraposisi mendaratkan dua pukulan cambukan secara serempak maut: Ia membalik mutasi kursi (seperti Konvers) \textbf{DAN} serentak menusuknya peluru Negasi (seperti Invers).

\begin{contoh}[Kontraposisi Sang Saudara Kembar Validitas]
Bermodal Implikasi Hujan-Aspal Thamrin di atas ($p \rightarrow q$): \\
\textbf{Varian Kontraposisi ($\neg q \rightarrow \neg p$):} "Oh? Kau menemukan fakta aspal mawar Thamrin detik ini tampak gersang kerontang sama sekali Tidak Basah ($\neg q$)? Hantam putus asaku, detik ini mutlak \textbf{MUSTAHIL} dewa langit semesta sedang iseng mengguyurkan meraup lebat Hujan di atas sana ($\neg p$)!!"

Adil dan agung! Kesimpulannya 100\% sah mutlak.
\textbf{Kesimpulan Aljabar Puncaknya:} Seumur hidup jagat logika komputasi, EKUIVALENSI ($\equiv$) mutlak cuma terjalin mengikat relasi cinta darah daging sehati sedarah kembar siam antara $p \rightarrow q$ murni dengan bayangan Kontraposisi $\neg q \rightarrow \neg p$. 
\end{contoh}

\subsection{Hukum Penyangkal Majemuk (Negasi Menembus Kurung)}

Untuk meruntuhkan tembok benteng komposit panjang yang dilapisi tanda Kurung di alam Aljabar Logika Hukum De Morgan senantiasa dipanggil ke mimbar, sedangkan pengingkaran gerbong panah Implikasi punya ritual sakral tersendiri yang murni anomali aneh.

\begin{itemize}
    \item \textbf{Negasi membongkar Konjungsi/Disjungsi (De Morgan Biasa)}
    \begin{align*}
        \neg(p \wedge q) &\equiv \neg p \vee \neg q \\
        \neg(p \vee q) &\equiv \neg p \wedge \neg q
    \end{align*}
    
    \item \textbf{Negasi membongkar perut panah Implikasi (Penting!)}
    Ingat rahasia di bab selanjutnya: Implikasi sejatinya adalah topeng disjungsi $(p \rightarrow q \equiv \neg p \lor q)$. Bila panah topeng raksasa ini ditembak mortir Negasi luar, ia bermetamorfos murni tumpul patah tak berdayung:
    \begin{align*}
        \neg(p \rightarrow q) &\equiv \neg(\neg p \vee q) \\
                              &\equiv p \wedge \neg q \quad \text{\textbf{(Hafalkan di Luar Kepala!)}}
    \end{align*}
    "Mengingkari janji Caleg *Jika Menang (p) Maka ia membagi Laptop (q)*, dinobatkan resmi patah berdosa bilamana \textbf{Ia sungguh Menang di mimbar (p)} DAN namun beralasan \textbf{Tak melempar Laptop Sebiji Pun ($\neg q$)}."
\end{itemize}
